\appendix

\section{The vacuum-polarization sign, from the quadratic effective action}
\label{app:sign}

The sign is derived once and propagated; it is not adjusted to taste. A unit
test (\nolinkurl{tests/test_sign_convention.py}) fails if any step below is
changed inconsistently, and in particular asserts that the opposite convention
produces a \emph{negative} spectral measure and antiscreening.

Integrating out the matter sector gives, in Euclidean signature and on the
transverse subspace,
\begin{equation}
S^{(2)}=\frac12\int_{q}A_{\mu}(-q)
\Bigl[\tfrac{1}{g^{2}}\bigl(q^{2}\delta_{\mu\nu}-q_{\mu}q_{\nu}\bigr)
+\Pi_{\mu\nu}(q)\Bigr]A_{\nu}(q),
\qquad
\Pi_{\mu\nu}=\bigl(q^{2}\delta_{\mu\nu}-q_{\mu}q_{\nu}\bigr)\Pi_{E}(q^{2}).
\end{equation}
The transverse kernel is $K_{T}(Q^{2})=Q^{2}\bigl(g^{-2}+\Pi_{E}(Q^{2})\bigr)$
and the static response is its inverse,
\begin{equation}
\Gcal(Q^{2})=\frac{g^{2}}{Q^{2}\bigl[1+g^{2}\Pi_{E}(Q^{2})\bigr]}.
\label{eq:Kinv}
\end{equation}
For the ungauged matter current, fix
\begin{equation}
i\!\int\! d^{4}x\,e^{ipx}\langle T j_{\mu}(x)j_{\nu}(0)\rangle
=(p_{\mu}p_{\nu}-p^{2}\eta_{\mu\nu})\Pi(p^{2}),
\qquad
\Pi(p^{2})-\Pi(0)=p^{2}\!\int\!\frac{\rho_{J}(s)\,ds}{s(s-p^{2}-i0)},
\end{equation}
with $\rho_{J}\ge0$ by Lehmann positivity, assuming one subtraction suffices.
Continuing to spacelike $p^{2}=-Q^{2}$ and defining
\begin{equation}
\overline{\Pi}(Q^{2})\equiv-\bigl[\Pi(-Q^{2})-\Pi(0)\bigr]
=Q^{2}\int\frac{\rho_{J}(s)\,ds}{s(s+Q^{2})}\ \ge\ 0,
\end{equation}
which is non-negative and increasing, and renormalising $\Pi_{E}(0)=0$ so that
$g=g_{R}$ is the long-distance coupling, \eqref{eq:Kinv} becomes
\begin{equation}
\Gcal(Q^{2})=\frac{g_{R}^{2}}{Q^{2}\bigl[1-g_{R}^{2}\overline{\Pi}(Q^{2})\bigr]}
=\frac{g_{R}^{2}}{Q^{2}}
+g_{R}^{4}\int\frac{\rho_{J}(s)\,ds}{s(s+Q^{2})}+O(g_{R}^{6}),
\label{eq:expansion}
\end{equation}
so that
\begin{equation}
\dsig_{\mathrm{LO}}(s)=g_{R}^{4}\,\frac{\rho_{J}(s)}{s}\,ds\ \ge\ 0 ,
\end{equation}
which is Hypothesis~\ref{hyp:stieltjes} at leading order. The next order does
not inherit this. The $O(g_R^6)$ term of~\eqref{eq:expansion} is
$g_R^6Q^2\bigl[\int\rho_J\,ds/(s(s+Q^2))\bigr]^2$, which vanishes at $Q^2=0$,
rises and then falls like $1/Q^2$. A function $\int\dsig/(Q^2+s)$ with
$\dsig\ge0$ is largest at $Q^2=0$, so this term is not of that form, and
truncating~\eqref{eq:expansion} at fixed order does not preserve the
hypothesis. Whatever positivity the chain has belongs to its sum.

\paragraph{The resummed bubble chain.}
Writing $Q^2/[s(s+Q^2)]=1/s-1/(s+Q^2)$ puts the denominator
of~\eqref{eq:expansion} in the form
\begin{equation}
W(Q^2)\equiv1-g_R^2\overline{\Pi}(Q^2)
=Z_3+g_R^2\int\frac{\rho_J(s)\,ds}{s+Q^2},
\qquad
Z_3\equiv1-g_R^2\int\frac{\rho_J(s)}{s}\,ds ,
\label{eq:Wpf}
\end{equation}
so that $W(0)=1$, $W(Q^2)\to Z_3$ as $Q^2\to\infty$, and $W$ decreases
strictly on $Q^2>0$ whenever $\rho_J\ge0$ does not vanish identically.
The sign of $Z_3$ decides the outcome.

If $Z_3>0$, $W/g_R^2$ is a Stieltjes function. By the duality between Stieltjes
and complete Bernstein functions~\cite{schilling_song_vondracek2012},
$Q^2W/g_R^2=1/\Gcal$ is then complete Bernstein and $\Gcal$ is Stieltjes. Its
Coulomb residue is $g_R^2/W(0)=g_R^2$, and $0<\Gcal\le g_R^2/(Q^2Z_3)$ excludes
an additive constant, so the resummed chain satisfies
Hypothesis~\ref{hyp:stieltjes} to all orders in the chain.

If $Z_3<0$, $W$ has exactly one zero, at some $Q^2_{\mathrm g}>0$. There $\Gcal$
has a simple pole with residue $g_R^2/[Q^2_{\mathrm g}W'(Q^2_{\mathrm g})]<0$, at
spacelike momentum, which no Stieltjes function admits.

If $Z_3=0$ there is no such pole, but
$\Gcal(Q^2)=\bigl[\int Q^2\rho_J\,ds/(s+Q^2)\bigr]^{-1}\to1/\int\rho_J\,ds$,
a positive constant whenever that total mass is finite, and the hypothesis
excludes a constant. The absence of a spacelike pole is therefore necessary for
Hypothesis~\ref{hyp:stieltjes} and not sufficient.

Two consequences bear on the rest of the paper. The continuum spectral density
of the chain is $g_R^4\rho_J(s)/(s\,|W(-s-i0)|^2)$, non-negative for every
coupling, $Z_3<0$ included. A positivity test built from moments of the
continuum therefore cannot see the failure at $Z_3<0$, which lives in a discrete
pole; the finite-moment conditions of Section~\ref{sec:scope} have exactly this
limitation. Second, if $\rho_J$ is cut off sharply at $s=\Lambda^2$ and does not
vanish there, as the Dirac density does not, then $W$ is real again for
$Q^2<-\Lambda^2$ and, when $Z_3>0$, vanishes once in that range. The resummed
measure then carries an atom of positive weight above the cutoff, outside the
support of $\rho_J$. This is the mechanism studied
in~\cite{giacosa_wolkanowski_2012}, where resummation places a pole on the
physical sheet outside the input spectrum, with positive residue, and a
spectral sum rule restores the normalisation. Theorem~\ref{thm:A} already
admits atoms in $\sigma$ through~\eqref{eq:pushforward}.

For the one-loop Dirac density at $\alpha=1/137$, $Z_3=0.966$ at
$\Lambda^2=10^{20}\,m^2$, and $Z_3$ reaches zero only near
$\log(\Lambda^2/4m^2)=3\pi/\alpha$, the Landau pole of the one-loop coupling.
Below it the chain satisfies the hypothesis with a wide margin. None of this
bears on the irreducible $O(g_R^6)$ contributions or on the nonperturbative
kernel, where Hypothesis~\ref{hyp:stieltjes} remains an assumption.

\paragraph{Check 1 (inverse kernel).}
The $O(g_{R}^{4})$ coefficient in \eqref{eq:expansion} is
$+\rho_{J}/[s(s+Q^{2})]$. Writing $1+g_{R}^{2}\overline{\Pi}$ in place of
$1-g_{R}^{2}\overline{\Pi}$ --- with $\overline{\Pi}\ge0$ --- flips it to
$-\rho_{J}/[s(s+Q^{2})]$ and destroys the hypothesis. This is the
discriminating check.

\paragraph{Check 2 (Uehling).}
Inserting the one-loop Dirac density
$\rho_{J}=\frac{q^{2}}{12\pi^{2}}(1+2m^{2}/s)\sqrt{1-4m^{2}/s}$ and
substituting $s=4m^{2}\zeta^{2}$ gives the potential correction prefactor
$g_{R}^{2}q^{2}/6\pi^{2}$, which for $q=1$ and $g_{R}^{2}=4\pi\alpha$ equals
$2\alpha/3\pi$: the Uehling coefficient~\cite{uehling_1935}.

\paragraph{Check 3 (screening).}
$g^{2}_{\mathrm{eff}}(Q^{2})=Q^{2}\Gcal(Q^{2})=g_{R}^{2}/[1-g_{R}^{2}
\overline{\Pi}]$ is increasing in $Q^{2}$, so the coupling grows in the
ultraviolet. The flipped convention gives a decreasing effective coupling,
i.e. antiscreening --- the signature of the error.

\paragraph{One-species cap.}\label{app:cap}
For the one-loop Dirac density, put $u=4m^2/x^2\in[0,1]$ and
$f(u)=(1+u/2)\sqrt{1-u}$. Since
$f(u)^2=1-3u^2/4-u^3/4$, one has $0\le f(u)\le1$. Consequently
\begin{equation}
0\le\delta_D(r)=\frac{g_R^2q^2}{6\pi^2}r^2
\int_{2m}^{\infty}x e^{-rx}f(4m^2/x^2)\,dx
\le\frac{g_R^2q^2}{6\pi^2}r^2\int_0^{\infty}x e^{-rx}\,dx
=\frac{g_R^2q^2}{6\pi^2}.
\end{equation}
This bounds the one-loop coefficient, not the uncomputed interacting remainder.

\section{Proofs}
\label{app:proofs}

\paragraph{Theorem~\ref{thm:C}.}
Write $d\mu_{r}=e^{-rx}\dnu/\Phi(r)$, a probability measure. Then
$\Gamma=-\Phi'/\Phi=\langle x\rangle_{r}\ge\Mstar$ since $\nu$ is supported in
$[\Mstar,\infty)$, and differentiating the ratio of moments gives
$\Gamma'=-(\langle x^{2}\rangle_{r}-\langle x\rangle_{r}^{2})=-\Var_{r}(x)$.
For convergence, fix $\varepsilon>0$ and $r_{0}>0$. The interval
$I=[\Mstar,\Mstar+\varepsilon/4]$ has mass $c>0$ when $\Mstar=\inf\supp\nu$, so
$\Phi(r)\ge c\,e^{-r(\Mstar+\varepsilon/4)}$. In the numerator of
$\Gamma-\Mstar$ the contribution from $x\le\Mstar+\varepsilon/2$ is at most
$(\varepsilon/2)\Phi$, and the tail is bounded by
$A\,e^{-(r-r_{0})(\Mstar+\varepsilon/2)}$ with
$A=\int(x-\Mstar)e^{-r_{0}x}\dnu<\infty$ by Hypothesis~\ref{hyp:growth}. The
ratio of the tail to $\Phi$ is $O(e^{-r\varepsilon/4})\to0$, so
$\limsup_{r}(\Gamma-\Mstar)\le\varepsilon/2$; letting $\varepsilon\to0$ and
using monotonicity gives the limit. \qed

\paragraph{Theorem~\ref{thm:E}, bound and monotonicity.}
For $p_{v}(x)=\sum_{i\le K}v_{i}x^{i}$,
\begin{equation}
v^{\mathsf{T}}(H_{1}-\Mstar H_{0})v
=\int(x-\Mstar)\,p_{v}(x)^{2}e^{-rx}\dnu\ \ge\ 0,
\end{equation}
since $x\ge\Mstar$ on the support; the Rayleigh characterisation and the bound
follow. Enlarging the polynomial space from degree $K$ to $K+1$ can only lower
an infimum, giving $B_{K+1}\le B_{K}$; and $K=0$ gives $a_{1}/a_{0}=\Gamma$.
$H_{0}\succ0$ provided the support contains at least $K+1$ points; for a
measure with $L$ atoms the pencil is exact at order $L-1$ and singular
beyond, which must be handled by rank reduction rather than reported as a new
bound. \qed

\paragraph{Theorem~\ref{thm:E}, convergence.}
Fix $r>0$. First, \emph{determinacy}: for $0<t<r$, using
$e^{tx}\ge(tx)^{2n}/(2n)!$,
\begin{equation}
a_{2n}(r)\le(2n)!\,t^{-2n}\,\Phi(r-t),
\qquad\text{hence}\qquad
\sum_{n\ge1}a_{2n}(r)^{-1/(2n)}=\infty,
\end{equation}
which is Carleman's criterion~\cite{akhiezer1965moment}. Carleman establishes
determinacy and nothing more; it is not by itself a proof of edge convergence.

Second, \emph{edge convergence}. Let $d\omega=(1+x)e^{-rx}\dnu$, which is
finite and possesses an exponential moment, $\int e^{tx}d\omega<\infty$ for
$0<t<r$. Polynomials are therefore dense in $L^{2}(\omega)$: if $f\perp$ all
polynomials, Cauchy--Schwarz makes $z\mapsto\int f(x)e^{zx}d\omega$ analytic in
a strip about the imaginary axis with all derivatives vanishing at the origin,
so it vanishes identically and $f=0$ by uniqueness of the Fourier transform.
Now let $\varepsilon>0$ and approximate in $L^{2}(\omega)$ the indicator of
$[\Mstar,\Mstar+\varepsilon]$, which has positive mass when
$\Mstar\in\supp\nu$. Because $d\omega$ dominates both $e^{-rx}\dnu$ and
$x\,e^{-rx}\dnu$, the $L^{2}(\omega)$ norm controls both integrals in the
Rayleigh quotient, so the quotient of the approximating polynomials converges
to that of the indicator, which is at most $\Mstar+\varepsilon$. Hence
$\lim_{K}B_{K}\le\Mstar+\varepsilon$ for every $\varepsilon>0$; with
$B_{K}\ge\Mstar$ the claim follows. Finally $\supp\mu_{r}=\supp\nu$ because
$e^{-rx}>0$ everywhere, so the edge reached is $\inf\supp\nu$. \qed

\paragraph{Edge law.}
Apply Watson's lemma~\cite{dlmf} separately to $\Phi(r)=\int e^{-rx}\dnu$ and
to $\int t\,e^{-rx}\dnu$ with $t=x-\Mstar$, then form
$\Gamma-\Mstar=(\int t\,e^{-rx}\dnu)/\Phi$. Subtracting $\Mstar\Phi$ before
expanding avoids differentiating an asymptotic series. With
$\dnu/dx=Ct^{p-1}[1+\beta t+\gamma t^{2}+O(t^{3})]$ and
$\int_{0}^{\infty}t^{p+j-1}e^{-rt}dt=\Gamma_{E}(p+j)r^{-p-j}$, dividing the two
series gives
\begin{equation}
\Gamma(r)=\Mstar+\frac{p}{r}+\frac{\beta p}{r^{2}}
+\frac{2\gamma p(p+1)-\beta^{2}p^{2}}{r^{3}}+O(r^{-4}).
\end{equation}
If only $1+O(t)$ is known, only $\Gamma=\Mstar+p/r+O(r^{-2})$ follows. For an
atom $Z\delta_{\Mstar}$ \emph{separated} from the rest of the support by
$d>0$, convergence is $O(e^{-dr})$; separation is a hypothesis, not a
consequence of the atom. If an atom and a continuum share the same edge,
$\Gamma-\Mstar\sim C\Gamma_{E}(p+1)r^{-p-1}/Z$: algebraic despite the pole.
\qed
